\begin{proof}[Proof of \cref{obj:thm:honest-lower-full-elbow}]
\begin{enumerate}
\item Choose the amplitude supplied by \cref{obj:lem:legal-iv-mixture-full-elbow}, and define
\[
0<c_{\star}<1,
\qquad
c_0=(1-\alpha)2^{3/4}c_{\star}^2>0.
\]
The amplitude is chosen uniformly in \(n\) and \(a\), so \(c_0\) depends on the prescribed constants \(\alpha,c_f,C_f,L\).
% lean: CausalSmith.Stat.TransportCaceRoughnuisanceLength.legal_iv_mixture_full_elbow

\item Fix \(n\ge256\) and \(0<a\le1/4\). Use the center, components, and observed mixtures of \cref{obj:lem:legal-iv-mixture-full-elbow}, with
\[
K_{\mathrm L}=\lceil n^{4/3}\rceil,
\qquad
\bar a=\max\{a,n^{-1/3}\},
\qquad
r_{\mathrm{sep}}=\frac{J_n}{\bar a}.
\]
Since the negative power is decreasing on the positive half-line,
\[
0<n^{-1/3}\le64^{-1/3}=\frac14,
\qquad
0<\bar a\le\frac14.
\]
The separation bound in that same result gives
\[
0<2^{3/4}c_{\star}^2n^{-1/3}\le J_n,
\qquad
r_{\mathrm{sep}}>0.
\]
For every sign vector \(\lambda\in\{-1,1\}^{K_{\mathrm L}}\), the component at \(\tau=-1\) belongs to \(\mathcal M_n(a)\subseteq\mathcal M_n\). Applying \cref{obj:lem:transported-cace-identification} to this component yields
\[
r_{\mathrm{sep}}
=\theta(Q_{\lambda,-1})
\in[-1,1].
\]
Consequently,
\[
0<r_{\mathrm{sep}}\le1,
\qquad
[-r_{\mathrm{sep}},r_{\mathrm{sep}}]\subseteq\Theta.
\]
The sampling conditions ensure that the center and component observed laws are probability laws.
% lean: CausalSmith.Stat.TransportCaceRoughnuisanceLength.actualStrength_bounds
% lean: CausalSmith.Stat.TransportCaceRoughnuisanceLength.transported_cace_identification

\item The strength floor implements the capped inverse-strength rate. Indeed, if \(a\le n^{-1/3}\), then
\[
\bar a=n^{-1/3},
\qquad
\frac{n^{-1/3}}{\bar a}=1
=\min\left\{1,\frac{n^{-1/3}}a\right\}.
\]
If \(a>n^{-1/3}\), then
\[
\bar a=a,
\qquad
\frac{n^{-1/3}}{\bar a}
=\frac{n^{-1/3}}a
=\min\left\{1,\frac{n^{-1/3}}a\right\}.
\]
Dividing the separation lower bound by \(\bar a>0\) and multiplying by \(1-\alpha>0\) therefore gives
\[
c_0\min\left\{1,\frac{n^{-1/3}}a\right\}
=(1-\alpha)2^{3/4}c_{\star}^2
  \frac{n^{-1/3}}{\bar a}
\le(1-\alpha)r_{\mathrm{sep}}.
\]
% lean: CausalSmith.Stat.TransportCaceRoughnuisanceLength.actualStrength_rate_identity

\item Let \(C_n\) be any procedure satisfying the theorem's hypotheses. For every \(t\in\mathbb R\), define the inclusion event
\[
E_t=\{\omega:t\in C_n(\omega)\}.
\]
These events are measurable by graph measurability. For \(t\in[-r_{\mathrm{sep}},r_{\mathrm{sep}}]\), define
\[
\tau_t=-\frac{t}{r_{\mathrm{sep}}}.
\]
The inequalities \(-r_{\mathrm{sep}}\le t\le r_{\mathrm{sep}}\), together with \(r_{\mathrm{sep}}>0\), give
\[
-1\le\tau_t\le1.
\]
By \cref{obj:lem:legal-iv-mixture-full-elbow}, every component indexed by \(\tau_t\) belongs to \(\mathcal M_n(a)\) and satisfies
\[
\theta(Q_{\lambda,\tau_t})
=-\frac{\tau_tJ_n}{\bar a}
=-\tau_t r_{\mathrm{sep}}
=t.
\]
Uniform honesty thus implies, for every sign vector,
\[
\left(
Q_{\lambda,\tau_t,\mathrm S}^{\otimes n}
\otimes Q_{\lambda,\tau_t,\mathrm T}^{\otimes n}
\right)(E_t)\ge1-\alpha.
\]
Averaging these inequalities over the finite sign family gives
\[
M_{\tau_t}(E_t)
=
2^{-K_{\mathrm L}}
\sum_{\lambda\in\{-1,1\}^{K_{\mathrm L}}}
\left(
Q_{\lambda,\tau_t,\mathrm S}^{\otimes n}
\otimes Q_{\lambda,\tau_t,\mathrm T}^{\otimes n}
\right)(E_t)
\ge1-\alpha.
\]
The normalized mixture is a probability law. The total-variation event inequality and the bound in \cref{obj:lem:legal-iv-mixture-full-elbow} yield
\[
M_{\tau_t}(E_t)-P_{\star}^{(n,n)}(E_t)
\le
\operatorname{TV}(M_{\tau_t},P_{\star}^{(n,n)})
\le\frac{1-\alpha}{2}.
\]
Hence, throughout the symmetric interval,
\[
P_{\star}^{(n,n)}(t\in C_n)
\ge\frac{1-\alpha}{2},
\qquad
t\in[-r_{\mathrm{sep}},r_{\mathrm{sep}}].
\]
% lean: CausalSmith.Stat.TransportCaceRoughnuisanceLength.uniform_data_mixture_coverage
% lean: Causalean.Stat.measureReal_sub_le_tvDist
% lean: CausalSmith.Stat.TransportCaceRoughnuisanceLength.honest_lower_full_elbow

\item Define the center's inclusion-probability function by
\[
p_{\star,n}(t)=P_{\star}^{(n,n)}(t\in C_n),
\qquad t\in\mathbb R.
\]
Graph measurability makes this function measurable, and
\[
0\le p_{\star,n}(t)\le1.
\]
It is therefore integrable over \(\Theta\). Applying \cref{obj:lem:random-set-length-identity} and using the interval inclusion established above gives
\[
\begin{aligned}
(1-\alpha)r_{\mathrm{sep}}
&=\int_{-r_{\mathrm{sep}}}^{r_{\mathrm{sep}}}
       \frac{1-\alpha}{2}\,dt\\
&\le\int_{-r_{\mathrm{sep}}}^{r_{\mathrm{sep}}}
       p_{\star,n}(t)\,dt\\
&\le\int_{\Theta}p_{\star,n}(t)\,dt\\
&=\mathbb E_{P_{\star}}
       \!\left[\lambda_{\Theta}(C_n)\right].
\end{aligned}
\]
The first inequality uses the pointwise inclusion bound; the second uses nonnegativity and
\([-r_{\mathrm{sep}},r_{\mathrm{sep}}]\subseteq\Theta\).
% lean: CausalSmith.Stat.TransportCaceRoughnuisanceLength.expectedLength_lower_of_inclusion

\item The range of expected lengths on the strength slice is bounded above. Pointwise,
\[
0\le\lambda_{\Theta}(C_n(\omega))
\le\operatorname{Leb}(\Theta)=2.
\]
For an integrable length, integration under an observed probability law gives
\[
\mathbb E_P\!\left[\lambda_{\Theta}(C_n)\right]
\le\int 2\,dP^{(n,n)}=2,
\qquad P\in\mathcal M_n(a).
\]
Under the convention assigning zero to a nonintegrable real integral, the complementary case also gives
\[
\mathbb E_P\!\left[\lambda_{\Theta}(C_n)\right]=0\le2.
\]
The measurable bounded lengths considered here are integrable. Since \(P_{\star}\in\mathcal M_n(a)\), the slice is nonempty, and comparison with its supremum now yields
\[
\begin{aligned}
c_0\min\left\{1,\frac{n^{-1/3}}a\right\}
&\le(1-\alpha)r_{\mathrm{sep}}\\
&\le\mathbb E_{P_{\star}}
       \!\left[\lambda_{\Theta}(C_n)\right]\\
&\le\sup_{P\in\mathcal M_n(a)}
       \mathbb E_P\!\left[\lambda_{\Theta}(C_n)\right]
\le2.
\end{aligned}
\]
This proves the asserted procedure-wise lower bound.
% lean: CausalSmith.Stat.TransportCaceRoughnuisanceLength.expectedLength_le_two
% lean: CausalSmith.Stat.TransportCaceRoughnuisanceLength.honest_lower_full_elbow

\item The honest interval class in \cref{obj:def:honest-intervals} is nonempty. Define the constant procedure by
\[
C_{n,\Theta}(\omega)=\Theta
\qquad\text{for every sample realization }\omega.
\]
It has a measurable graph and connected values. By \cref{obj:lem:transported-cace-identification},
\[
\theta(P)\in\Theta,
\qquad
P^{(n,n)}\bigl(\theta(P)\in C_{n,\Theta}\bigr)
=1\ge1-\alpha,
\qquad P\in\mathcal M_n.
\]
Thus
\[
C_{n,\Theta}\in\mathfrak H_n,
\qquad
\mathfrak H_n\ne\varnothing.
\]
Every member of this class satisfies the procedure-wise bound already established. Taking the infimum prescribed by \cref{obj:def:length-frontier} gives
\[
L_n(a)
=
\inf_{C_n\in\mathfrak H_n}
\sup_{P\in\mathcal M_n(a)}
\mathbb E_P\!\left[\lambda_{\Theta}(C_n)\right]
\ge
c_0\min\left\{1,\frac{n^{-1/3}}a\right\}.
\]
% lean: CausalSmith.Stat.TransportCaceRoughnuisanceLength.parameterSpace_honest
% lean: CausalSmith.Stat.TransportCaceRoughnuisanceLength.honest_lower_full_elbow
\end{enumerate}
\end{proof}
